\chapter{Panduan Karir AI: Roadmap \& Strategi}
\label{chap:career}

Karir di bidang Artificial Intelligence (AI) sangat menjanjikan, namun seringkali membingungkan karena banyaknya istilah dan peran yang tumpang tindih. Panduan ini dirancang untuk memberikan peta jalan (roadmap) yang jelas bagi Anda yang ingin berkarir di industri ini, baik sebagai pemula maupun profesional yang ingin beralih karir.

\section{Peta Peran (Role Map) dalam Industri AI}

Secara umum, pekerjaan di bidang AI dapat dibagi menjadi tiga kategori utama berdasarkan fokus utamanya: Analisis, Rekayasa (Engineering), dan Penelitian.

\subsection{1. Data Analyst}
\textbf{Fokus:} Menerjemahkan data menjadi wawasan bisnis (business insights).
\begin{itemize}
    \item \textbf{Tugas Utama:} Membersihkan data, membuat dashboard visualisasi (Tableau/PowerBI), melakukan analisis statistik deskriptif, dan mempresentasikan temuan kepada stakeholder non-teknis.
    \item \textbf{Skill Kunci:} SQL (wajib), Excel/Spreadsheet, Visualisasi Data, Python/R dasar, Statistik Dasar, Komunikasi Bisnis.
    \item \textbf{Cocok untuk:} Mereka yang suka bercerita dengan data dan memiliki ketajaman bisnis yang kuat.
\end{itemize}

\subsection{2. Data Scientist}
\textbf{Fokus:} Membangun model prediktif dan eksperimen statistik untuk memecahkan masalah kompleks.
\begin{itemize}
    \item \textbf{Tugas Utama:} Melakukan analisis eksploratif mendalam, membangun model Machine Learning (klasifikasi, regresi, clustering), melakukan A/B testing, dan mengoptimalkan metrik bisnis.
    \item \textbf{Skill Kunci:} Python (Pandas, Scikit-Learn), SQL Lanjut, Matematika (Aljabar Linear, Kalkulus, Probabilitas), Statistik Inferensial, Algoritma ML.
    \item \textbf{Cocok untuk:} Pemecah masalah yang memiliki latar belakang kuat di matematika/statistik dan ingin menerapkan metode ilmiah pada data.
\end{itemize}

\subsection{3. Machine Learning Engineer (MLE)}
\textbf{Fokus:} Mengambil model eksperimental dari Data Scientist dan membuatnya siap produksi (scalable, reliable, efficient).
\begin{itemize}
    \item \textbf{Tugas Utama:} Membangun pipeline data (ETL), men-deploy model ke cloud (AWS/GCP), memantau kinerja model di produksi (MLOps), mengoptimalkan inferensi (latency/throughput).
    \item \textbf{Skill Kunci:} Python Lanjut, Software Engineering (OOP, Design Patterns), Cloud Computing (Docker, Kubernetes), Framework ML (TensorFlow/PyTorch), CI/CD.
    \item \textbf{Cocok untuk:} Software Engineer yang tertarik dengan data, atau Data Scientist yang ingin fokus pada sistem dan infrastruktur.
\end{itemize}

\subsection{4. AI Research Scientist}
\textbf{Fokus:} Menemukan algoritma baru dan mendorong batas pengetahuan di bidang AI.
\begin{itemize}
    \item \textbf{Tugas Utama:} Membaca dan menulis paper akademis, merancang arsitektur model baru, bereksperimen dengan teknik pelatihan novel.
    \item \textbf{Skill Kunci:} Matematika Tingkat Lanjut, Pemahaman mendalam teori Deep Learning, Kemampuan implementasi paper (PyTorch/JAX), Riset Akademis.
    \item \textbf{Kualifikasi:} Seringkali membutuhkan gelar PhD atau Master di bidang terkait.
\end{itemize}

\subsection{5. AI Product Manager}
\textbf{Fokus:} Menjembatani kesenjangan antara tim teknis AI dan kebutuhan pengguna/bisnis.
\begin{itemize}
    \item \textbf{Tugas Utama:} Mendefinisikan roadmap produk AI, menentukan metrik keberhasilan, memahami etika AI, dan mengelola ekspektasi stakeholder.
    \item \textbf{Skill Kunci:} Pemahaman konsep AI (tanpa harus coding mendalam), Manajemen Produk, UX Design, Strategi Bisnis.
\end{itemize}

\section{Roadmap Belajar (Learning Path)}

Berikut adalah urutan belajar yang disarankan untuk menjadi ML Engineer atau Data Scientist dari nol:

\subsection{Tahap 1: Fondasi (Bulan 1-2)}
\begin{enumerate}
    \item \textbf{Python:} Sintaks dasar, fungsi, OOP, error handling.
    \item \textbf{Matematika:} Aljabar Linear (vektor, matriks), Kalkulus (turunan), Statistik \& Probabilitas dasar.
    \item \textbf{Data Manipulation:} NumPy (array), Pandas (DataFrame).
    \item \textbf{Visualisasi:} Matplotlib, Seaborn.
\end{enumerate}

\subsection{Tahap 2: Machine Learning Klasik (Bulan 3-4)}
\begin{enumerate}
    \item \textbf{Supervised Learning:} Regresi Linear, Logistik, Decision Trees, KNN, SVM.
    \item \textbf{Unsupervised Learning:} K-Means Clustering, PCA.
    \item \textbf{Evaluasi Model:} Train/Test Split, Cross-Validation, Precision/Recall, ROC-AUC.
    \item \textbf{Project:} Prediksi harga rumah (Regresi), Klasifikasi Iris/Titanic.
\end{enumerate}

\subsection{Tahap 3: Deep Learning \& Neural Networks (Bulan 5-6)}
\begin{enumerate}
    \item \textbf{Teori:} Perceptron, Backpropagation, Activation Functions, Loss Functions.
    \item \textbf{Framework:} PyTorch (disarankan untuk riset/belajar) atau TensorFlow/Keras (untuk industri tertentu).
    \item \textbf{Arsitektur:} CNN (untuk gambar), RNN/LSTM (untuk teks/waktu).
    \item \textbf{Project:} Klasifikasi gambar (CIFAR-10), Sentiment Analysis.
\end{enumerate}

\subsection{Tahap 4: Spesialisasi \& Modern AI (Bulan 7+)}
Pilih satu jalur spesialisasi:
\begin{itemize}
    \item \textbf{NLP (Natural Language Processing):} Transformers (BERT, GPT), LLM, Hugging Face, LangChain.
    \item \textbf{Computer Vision:} Object Detection (YOLO), Segmentation, Generative Models (GAN, Diffusion).
    \item \textbf{Reinforcement Learning:} Q-Learning, Policy Gradient.
\end{itemize}

\subsection{Tahap 5: MLOps \& Engineering (Parallel)}
Pelajari cara kerja sistem produksi:
\begin{enumerate}
    \item \textbf{Git:} Version control wajib.
    \item \textbf{API Development:} Flask atau FastAPI.
    \item \textbf{Docker:} Containerization.
    \item \textbf{Cloud:} AWS (S3, EC2, SageMaker) atau GCP.
\end{enumerate}

\section{Membangun Portofolio}

Gelar akademik penting, tapi portofolio proyek nyata seringkali lebih meyakinkan bagi rekruter.
\begin{itemize}
    \item \textbf{Kualitas > Kuantitas:} Lebih baik punya 2-3 proyek end-to-end yang solid daripada 10 notebook tutorial yang belum selesai.
    \item \textbf{End-to-End:} Jangan berhenti di file \texttt{.ipynb}. Buat aplikasi web sederhana (Streamlit) agar orang bisa mencoba model Anda. Deploy ke cloud (gratisan seperti Hugging Face Spaces atau Streamlit Cloud).
    \item \textbf{Dokumentasi:} Tulis README di GitHub dengan jelas: Masalah apa yang diselesaikan? Data dari mana? Bagaimana cara menjalankannya? Apa hasilnya?
    \item \textbf{Blog Post:} Tulis artikel Medium atau LinkedIn yang menjelaskan perjalanan Anda memecahkan masalah tersebut.
\end{itemize}

\section{Persiapan Wawancara}

\subsection{Coding Interview (LeetCode)}
Banyak perusahaan teknologi besar masih menggunakan tes algoritma standar. Fokus pada struktur data (Array, String, Hash Map, Tree) dan algoritma dasar. Python adalah bahasa terbaik untuk ini.

\subsection{Machine Learning Design Interview}
Anda akan diberi masalah abstrak (misal: "Desain sistem rekomendasi berita untuk homepage"). Kerangka jawabannya:
1. \textbf{Klarifikasi Masalah:} Metrik bisnis apa? Latency requirement? Ukuran data?
2. \textbf{Data:} Sumber data? Labeling strategy? Feature engineering?
3. \textbf{Model:} Baseline model? Model kompleks? Mengapa memilih itu?
4. \textbf{Evaluasi:} Metrik offline (AUC, RMSE) vs online (CTR, Watch Time).
5. \textbf{Deployment:} Batch vs Real-time inference? Monitoring drift?

\section{Gaji \& Prospek Pasar (Indonesia \& Global)}
\textit{Catatan: Angka ini adalah estimasi kasar per 2024 dan sangat bervariasi tergantung lokasi dan ukuran perusahaan.}

\subsection{Indonesia (Per Bulan)}
\begin{itemize}
    \item \textbf{Junior (0-2 tahun):} Rp 8.000.000 - Rp 15.000.000
    \item \textbf{Mid-Level (2-5 tahun):} Rp 15.000.000 - Rp 30.000.000
    \item \textbf{Senior (5+ tahun):} Rp 30.000.000 - Rp 60.000.000+
\end{itemize}

\subsection{Global/Remote (Per Tahun - USD)}
\begin{itemize}
    \item \textbf{Junior:} \$60.000 - \$100.000
    \item \textbf{Senior:} \$120.000 - \$200.000+
    \item \textbf{AI Researcher (Top Tech):} \$300.000+
\end{itemize}

\section{Sumber Belajar Rekomendasi}
\begin{itemize}
    \item \textbf{Kursus:} Coursera (Andrew Ng), Fast.ai (Jeremy Howard), Udacity.
    \item \textbf{Buku:} "Hands-On Machine Learning" (Aurelien Geron), "Deep Learning with Python" (Francois Chollet).
    \item \textbf{Komunitas:} Kaggle, Reddit r/MachineLearning, Discord server (OpenAI, Hugging Face).
\end{itemize}

\begin{tipbox}
    \textbf{Saran Terakhir:} Jangan menunggu sampai merasa "siap 100\%" untuk melamar. Bidang ini terlalu luas untuk dikuasai sepenuhnya. Mulailah membangun, bagikan karya Anda, dan belajar sambil berjalan.
\end{tipbox}
